\chapter{Analytical proposition and evaluation questions}
\label{cap:hipotesis}

The study distinguishes an algebraic consequence of the swap rule from two questions that require simulation. These are claims about the specified zero-fee micro-world; they are not hypotheses about observed markets.

\section{H1: ideal correction under bounded arbitrage}

Let $e_t=\log(p_t^D/p_t^C)$ and $D_t=(|e_t|-\gamma)_+$. In an isolated, zero-fee pool with a fixed reference, no other orders and no binding inventory limit, the maximum order $q=cx$ changes the log price by $\delta_c=2\log(1+c)$. Since the arbitrageur stops at the band, the out-of-band error evolves as
\[
  D_{t+1}=\max(0,D_t-\delta_c),\qquad
  N_\gamma=\left\lceil\frac{D_0}{2\log(1+c)}\right\rceil \quad(c>0).
\]
For $c=0$, no finite correction occurs. This is a mechanical benchmark derived from the constant-product swap equation, not an empirical finding. Simulation is useful only for quantifying departures from this benchmark caused by other agents, event order, finite horizon and the specified inventory rules.

\section{H2: withdrawal, order scale and tracking response}

The existing E3 rule is a one-time threshold trigger, not an estimated adaptive policy: it withdraws 45\,\% of reserves as soon as $|e_t|>0.02$. The executed shocks exceed this threshold in every E3 run. The question addressed by the available runs is therefore narrower: how does this fixed withdrawal rule change the simulated response under fixed absolute order sizes, and how does that contrast change when order sizes scale with reserves or the panic action is disabled? The original peak displacement from the pre-shock AMM price is treated as a diagnostic only; it is not a measure of overshoot relative to the new reference.

\section{H3: predictive value of reduced maps}

The held-out comparison asks whether the fitted two-state map improves one-step prediction over persistence and a scalar AR(1), and whether band-dependent terms improve recursive replay error over an affine map. This comparison concerns simulator-generated paths. The hybrid replay receives the realized future CEX path, so its result is not a real-time forecast test. No regime-classification or stability hypothesis is advanced.
